\documentclass[margin,line]{res}
\usepackage[T1]{fontenc}
\usepackage{lmodern}
\usepackage[hidelinks]{hyperref}
\oddsidemargin -.5in
\evensidemargin -.5in
\textwidth=6.0in
\itemsep=0in
\parsep=0in

% Consistent spacing between entries
\newcommand{\entrygap}{\vspace{.09in}}

% Paper entry: {title}{year}{authors line}
\newcommand{\paper}[3]{%
  {\em ``#1''} \hfill {\bf #2}\\
  {\small #3}\\}

% Position entry: {title}{dates}
\newcommand{\role}[2]{{\em #1} \hfill {\bf #2}\\}

\begin{document}

\name{Blaize Giangiulio \vspace*{.1in}}

\begin{resume}

%-----------------------------------------------------------------
\section{\sc Contact\\ Information}
\vspace{.05in}
\begin{tabular}{@{}p{2.55in}p{2.6in}}
School of Economics                  & {\it E-mail:} \href{mailto:bg654@drexel.edu}{bg654@drexel.edu} \\
LeBow College of Business            & {\it Phone:} (856) 381-2348 \\
Drexel University, Office 1030       & {\it Website:} \href{https://bgian123.github.io}{bgian123.github.io} \\
3220 Market St, Philadelphia, PA 19104 \\
\end{tabular}

%-----------------------------------------------------------------
\section{\sc Education}
{\bf Drexel University}, Philadelphia, PA \hfill {\bf 2022--2027 (expected)}\\
Ph.D., Economics; Ph.D. Candidate since September 2024\\
Expected completion: June 2027\\
Advisor: Yoto V. Yotov

\entrygap
{\bf McGill University}, Montreal, QC, Canada\\
M.A., Economics \hfill {\bf 2021}\\
B.A., Economics, with Distinction \hfill {\bf 2020}

%-----------------------------------------------------------------
\section{\sc Fields}
International Trade, Applied Industrial Organization, Applied Econometrics

%-----------------------------------------------------------------
\section{\sc Job Market\\ Paper}
\paper{Chokepoints}{2026}{Solo-authored}
Models the disruption of a maritime chokepoint as a strategic choice rather than an exogenous shock. In an $(N+1)$-country Armington economy, a disruptor that raises trade costs between other countries improves its own terms of trade even absent any geopolitical motive, and the gain scales quadratically in the size of the disrupted corridor, rationalizing why disruptors target arteries rather than individual routes. Applies the model to the 2023--24 Red Sea crisis using a continuous, geography-determined measure of each country pair's exposure to closure of the Bab el-Mandeb, built by rerouting 39 million city-pair freight paths over a global road and maritime network. PPML gravity estimates show that trade fell on exposed dyads, with the response concentrated among Western-aligned pairs; exact hat algebra delivers the general-equilibrium welfare incidence, with losses increasing in exposure and Iran's own gain negligible. Inverting the disruptor's first-order condition shows that each dollar Iran gains at the margin costs Western economies roughly \$200 to \$1,200, and that self-interest alone rationalizes the campaign only if it cost Iran less than \$42 million per year, well below independent estimates of Iranian support to the Houthis: a positive weight on harming adversaries is required to explain the observed disruption.

%-----------------------------------------------------------------
\section{\sc Working\\ Papers}
\paper{Perceived Protection: Tariff Spillovers to Product Substitutes in the Used Car Market}{2026}{with Sebastien Bradley, Mian Dai, and Yoto V. Yotov \quad {\em CGPA Working Paper 2026-16}}
Studies the spillover effects of the 2025 ``Liberation Day'' automotive tariffs on the untaxed U.S. used-car market. Using over eight million weekly web-scraped listings and a VIN-based crosswalk identifying each vehicle's country of assembly, documents an immediate and persistent 0.2--0.6\% increase in the prices of foreign-brand used vehicles relative to domestic-brand vehicles, concentrated among the newest model years---the closest substitutes for new cars---and consistent with demand-side substitution. Foreign-brand vehicles assembled in the U.S. appreciate as much as those assembled abroad, suggesting that the perceived foreignness of a vehicle is largely independent of its true production location. Neglecting tariff pass-through to untaxed substitutes understates overall consumer incidence.

\entrygap
\paper{Logistic Distance}{2025--2026}{with Peter Herman, Serge Shikher, and Yoto V. Yotov}
Constructs a transportation-network-informed distance measure that incorporates realistic city-to-city routing into the Mayer and Zignago (2011) population-weighted distance formula, making bilateral distances responsive to infrastructure shocks. Applies the measure to the 2023 Red Sea crisis as a natural experiment, estimating a saturated gravity model in an event-study framework on monthly bilateral trade data.

\entrygap
\paper{Time-Varying Distance and Trade}{2025--2026}{with Peter Herman, Serge Shikher, and Yoto V. Yotov}
Extends the Mayer and Zignago (2011) population-weighted distance formula to allow bilateral distances to evolve over time using annual city-level population data from the mid-twentieth century onward, and examines within a structural gravity framework how long-run shifts in the spatial distribution of population have shaped bilateral trade patterns.

%-----------------------------------------------------------------
\section{\sc Work in\\ Progress}
\paper{Shrouded by Statute: Mandated Used-Car Warranties and the Market}{2026}{Solo-authored}
Asks whether state used-car lemon laws, which often require dealers to attach a short express warranty to vehicles sold below sharp odometer cutoffs (e.g., 125,000 miles in Massachusetts), generate price premiums. Using 9.9 million web-scraped listings and a regression-discontinuity design with rich fixed effects, finds no warranty premium, alongside pronounced bunching of odometer readings just above the Massachusetts cutoff and reallocation of near-cutoff inventory across state lines. Rationalizes both facts with a shrouded-attribute model in the spirit of Gabaix and Laibson (2006): when buyers unaware of the mandate are the marginal buyers of covered vehicles, the equilibrium premium is zero regardless of the dealer's cost; dealers have no private incentive to debias buyers because disclosure raises claims costs, cannibalizes service-contract sales, and is a public good among rivals; and the resulting unpriced cost acts as a tax on covered sales that is shifted onto the quantity margin. Calibrated simulations reproduce the observed magnitudes of both the null price effect and the Massachusetts bunching.

%-----------------------------------------------------------------
\section{\sc Presentations}
Philadelphia Trade Group, 2025\\
Drexel Sanctions Workshop (participant), 2024

%-----------------------------------------------------------------
\section{\sc Fellowships\\ and Awards}
Fellow, Center for Global Policy Analysis, Drexel University \hfill 2025--present\\
Dr.\ Tom Hindelang Outstanding Ph.D. Student Instructor Award, Drexel University \hfill 2025\\
Teaching Assistant Excellence Award, Drexel University \hfill 2024\\
Dragon Fellowship, Drexel University \hfill 2022--present\\
Master's Teaching Assistantship, McGill University \hfill 2020--2021

%-----------------------------------------------------------------
\section{\sc Teaching\\ Experience}
{\bf Instructor of Record}, Drexel University\\
Intermediate Microeconomics (ECON 301) \hfill Spring 2026\\
Principles of Microeconomics (ECON 201) \hfill Fall 2024, Fall 2025\\
Ph.D. Math Camp \hfill Summer 2025

\entrygap
{\bf Teaching Assistant}\\
Introductory Microeconomics, Drexel University \hfill 2022--2024\\
Introductory Macroeconomics, McGill University \hfill 2020--2021

%-----------------------------------------------------------------
\section{\sc Research\\ Experience}
{\bf Center for Global Policy Analysis}, Drexel University\\
\role{Fellow}{2025--present}
Contributed to the development of the large-scale EQUIP trade simulation model; developed a geospatially founded logistic distance measure; built distributed estimation pipelines on HPC clusters (Slurm), scaling to 5,000+ PPML regressions with high-dimensional fixed effects.

\entrygap
{\bf U.S. International Trade Commission}, Washington, DC\\
\role{Student Trainee (Economist)}{2024--2025}
Studied the effect of population movements on international trade flows in a structural gravity framework; used geospatial data and graph-theoretic methods in Python to predict shipping routes and construct a new logistic distance variable; provided technical assistance to USTR.

\entrygap
{\bf Drexel University}\\
\role{Research Assistant to Yoto V. Yotov}{2022--present}
Assembled, cleaned, and analyzed data from a wide range of trade sources in R and Stata; prepared summaries of trade-policy research.

%-----------------------------------------------------------------
\section{\sc Professional\\ Service}
Referee: {\em Defence and Peace Economics}, {\em Economics Letters}, {\em Review of International Economics}, {\em The World Economy}

%-----------------------------------------------------------------
\section{\sc Skills}
{\it Programming:} Stata, Python, R, MATLAB, Bash, Slurm (HPC), \LaTeX\\
{\it Languages:} English (native), French (fluent)

%-----------------------------------------------------------------
\section{\sc References}
\begin{tabular}{@{}p{3in}}
{\bf Yoto V. Yotov} (Advisor)\\
Professor of Economics\\
School of Economics, LeBow College of Business\\
Drexel University\\
\href{mailto:yotov@drexel.edu}{yotov@drexel.edu}\\
\end{tabular}

\end{resume}
\end{document}
